% GENERATED_BY: oc_core_monograph_translation_draft_pipeline_v1.py
% AI_ASSISTED_TRANSLATION_DRAFT: TRUE
% THEOREM_REVIEW_STATUS: PENDING
% THEOREM_REVIEW_MODE: FORMAL_AI_GATE__CODEX_CLI_CHATGPT
% THEOREM_REVIEW_REVIEWER_ID:
% THEOREM_REVIEW_AUTHORITY_CLASS:
% THEOREM_REVIEWED_AT:
% THEOREM_REVIEW_DOSSIER_REF:
% SOURCE_LANGUAGE: EN
% TARGET_LANGUAGE: RU
% SOURCE_EN_REF: content/k_levels/k1.tex
% ================================================================
% ==== FILE: content/k_levels/k1.tex
% ================================================================

% ==============================
%  Ontology of Continua — Core
%  K-level Module: K1
%  File: content/k_levels/k1.tex
%  Status: FULLY DEFINED
% ==============================

\subsubsection{\texorpdfstring{$K_1$}{K_1} Обзор}
\label{sec:k1-overview}

Уровень $K_1$ — это первый действительно
\emph{структурный} континуум.
Он возникает под действием оператора $\Psi_{0\to1}$ на $K_0$,
вводя одномерную ось $A_1$, тем самым позволяя существование
непрерывных степеней свободы.

$K_1$ ещё не является физическим пространством; скорее это абстрактный
прототип одномерного континуума:
\[
K_1 = (X, \tau, \mu, A_1),
\]
где $X$ — одномерная область, $\tau$ — структурная
топология, $\mu$ — мера, согласованная с $\tau$.

Создание $K_1$ — это первый шаг, когда:
\begin{itemize}
    \item различия могут быть выражены вдоль оси,
    \item возможно согласованное варьирование,
    \item структурное напряжение может распространяться,
    \item ограничения порождают непрерывные кривые.
\end{itemize}

Каждый более высокий континуум наследует $A_1$ как базовую ось.

% ================================================================
\subsubsection{\texorpdfstring{$\Omega(K_1)$}{\Omega(K_1)} Пространство состояний}

Пространство состояний $K_1$:
\[
\Omega(K_1) = C^0(X, V),
\]
то есть пространство непрерывных функций из одномерной области $X$
в множество значений $V$.

Свойства:
\begin{enumerate}
    \item $\Omega(K_1)$ бесконечномерно (пространство функций).
    \item По определению каждая конфигурация непрерывна.
    \item Непрерывности лежат вне $\Omega(K_1)$ и соответствуют
          нарушению структурной согласованности.
    \item $X$ может быть открытым, замкнутым, компактом или некомпактным
          в зависимости от ограничений соответствующего $M_1$.
\end{enumerate}

Это минимальное пространство, в котором существует непрерывное
варьирование.

% ================================================================
\subsubsection{\texorpdfstring{$\partial\Omega(K_1)$}{\partial\Omega(K_1)} Граница}

Граница состоит из конфигураций, приближающихся к потере непрерывности:
\[
\partial\Omega(K_1) =
\{ f \in C^0(X,V) : \exists\, x_0 \in X \\text{ при этом }
\lim_{x \to x_0} f(x) \\text{ не определено}\ \}.
\]

Интерпретация:
\begin{itemize}
    \item точки, где непрерывность нарушена;
    \item точки, где структурное напряжение превышает $\Theta_1$;
    \item точки, где область $X$ вырождается или схлопывается.
\end{itemize}

Переход границы уничтожает континуум $K_1$.

% ================================================================
\subsubsection{\texorpdfstring{$A(K_1)$}{A(K_1)} Оси}

У $K_1$ ровно одна ось:
\[
A_1 : X \to \mathbb{R},
\]
представляющая внутренний параметр, вдоль которого упорядочиваются
различия.

Условия:
\begin{itemize}
    \item $A_1$ должна быть монотонной (аксиома A$_{1,\text{mon}}$).
    \item $A_1$ должна быть связной (аксиома A$_{1,\text{conn}}$).
    \item $A_1$ должна допускать вариации, совместимые с $\tau$.
\end{itemize}

Все будущие оси $A_k$ для $k \ge 2$ расширяют эту базовую структуру.

% ================================================================
\subsubsection{\texorpdfstring{$P(K_1)$}{P(K_1)} Потенциалы}

Потенциалы на $K_1$ включают:
\[
P(K_1) = \{ P_{\mathrm{grad}}, P_{\mathrm{bound}}, P_{\mathrm{smooth}} \},
\]
описывая:
\begin{itemize}
    \item градиенты вдоль оси,
    \item напряжения, обусловленные границами,
    \item ограничения гладкости,
    \item допустимую кривизну функций в $\Omega(K_1)$.
\end{itemize}

Эти потенциалы ещё не силы; это структурные потенциалы, управляющие
согласованностью вариаций.

% ================================================================
\subsubsection{\texorpdfstring{$\Theta(K_1)$}{\Theta(K_1)} Пороги}

Пороги регулируют допустимость состояний:
\[
\Theta(K_1) = \{ \Theta_{\mathrm{cont}}, \Theta_{\mathrm{grad}},
                 \Theta_{\mathrm{smooth}} \}.
\]

Примеры:
\begin{itemize}
    \item $\Theta_{\mathrm{cont}}$ — минимальная непрерывность.
    \item $\Theta_{\mathrm{grad}}$ — максимальный допустимый градиент до
          того, как напряжение становится неустойчивым.
    \item $\Theta_{\mathrm{smooth}}$ — минимальная регулярность.
\end{itemize}

Если любой порог превышен, конфигурация оказывается вне $\Omega(K_1)$.

% ================================================================
\subsubsection{\texorpdfstring{$J(K_1)$}{J(K_1)} Потоки}

Потоки описывают допустимые непрерывные изменения вдоль $A_1$:
\[
J(K_1) = \{ \partial_x f(x),\; \partial_t f(x,t),\; \text{допустимые деформации} \}.
\]

Ключевые свойства:
\begin{itemize}
    \item потоки распространяют структурное напряжение,
    \item потоки должны сохранять непрерывность,
    \item потоки не могут создавать новые оси (требуется переход $K_1\to K_2$).
\end{itemize}

% ================================================================
\subsubsection{\texorpdfstring{$C(K_1)$}{C(K_1)} Циклы}

Поскольку $K_1$ поддерживает непрерывную вариацию, но ещё не многомерное
взаимодействие, циклы топологически тривиальны:
\[
C_{\mathrm{triv}} = \text{\text{постоянные петли в пространстве функций}}.
\]

Нет нетривиальных структурных циклов, поскольку циклическая структура
требует как минимум двух независимых осей ($K_2$).

Следовательно:
\[
\tau(K_1) \text{ не может возникать}.
\]

% ================================================================
\subsubsection{\texorpdfstring{$\tau(K_1)$}{\tau(K_1)} Время}

Время отсутствует в $K_1$.

Согласно теореме происхождения времени (Physics Run):

\[
\text{время возникает только на } K_2 \text{ из неразрушаемых циклов}.
\]

Следовательно:
\[
\tau(K_1) \text{ не определено и не представимо}.
\]

% ================================================================
\subsubsection{\texorpdfstring{$k(K_1)$}{k(K_1)} Континуумность}

Мера континуумности:
\[
k(K_1) =
H(\Omega(K_1)) \cdot
\frac{|A(K_1)|}{|A(M_1)|} \cdot
(1 - \frac{T_1}{\Theta_{\mathrm{cont}}})_+,
\]
где:
\begin{itemize}
    \item $|A(K_1)| = 1$,
    \item $A(M_1)$ может допускать больше осей, которые пока не реализованы,
    \item $T_1$ — структурное напряжение,
    \item $(\cdot)_+$ — положительная часть.
\end{itemize}

$k(K_1)$ измеряет, остаётся ли $K_1$ структурно жизнеспособным.

% ================================================================
\subsubsection{\texorpdfstring{$T(K_1)$}{T(K_1)} Структурное напряжение}

Структурное напряжение в $K_1$ возникает из:
\begin{itemize}
    \item градиентов функций,
    \item эффектов границ,
    \item несовместимости ограничений вдоль $A_1$.
\end{itemize}

Если:
\[
T(K_1) > \Theta_{\mathrm{cont}},
\]
то непрерывность теряется и континуум схлопывается.

% ================================================================
\subsubsection{\texorpdfstring{$E(K_1)$}{E(K_1)} Энергия}

Энергия определена структурно как:
\[
E(K_1) = \int_X \mathcal{E}(f(x), \partial_x f(x))\, d\mu,
\]
где $\mathcal{E}$ — плотность структурной энергии, зависящая от:
\begin{itemize}
    \item изменения значения,
    \item ограничений гладкости,
    \item штрафов регулярности.
\end{itemize}

$E(K_1)$ ещё не имеет физического смысла; это функционал, измеряющий
согласованность.

% ================================================================
\subsubsection{\texorpdfstring{$K_1$}{K_1} Операторы (\texorpdfstring{$\Psi$}{\Psi}, \texorpdfstring{$\Phi$}{\Phi}, \texorpdfstring{$\Lambda$}{\Lambda}, \texorpdfstring{$U$}{U}, \texorpdfstring{$\Chi$}{\Chi})}

\begin{itemize}
    \item $\Psi_{1\to2}$ — рождение второй оси $A_2$, позволяющей
          $K_2$.
    \item $\Phi$ — эволюция допустимого мета-пространства $M_1$.
    \item $\Lambda$ — оператор ограничения, обеспечивающий непрерывность,
          монотонность и связность.
    \item $U$ — унификация локальной структуры в глобальную согласованность.
    \item $\Chi$ — оператор ветвления, обеспечивающий несовместимые
          многообразия $K_1$.
\end{itemize}

Эти операторы определяют, как $K_1$ эволюционирует и как возникают новые
размерности.

% ================================================================
\subsubsection{Процессы на \texorpdfstring{$K_1$}{K_1}}

Ключевые процессы:
\begin{itemize}
    \item распространение градиентов,
    \item сглаживание разрывов,
    \item стабилизация непрерывных конфигураций,
    \item структурный коллапс при чрезмерном напряжении,
    \item переход на размерность $K_2$ под давлением
          несовместимых различий (аксиома~15).
\end{itemize}

% ================================================================
\subsubsection{Прогнозы \texorpdfstring{$K_1$}{K_1}}

Теория предсказывает:
\begin{enumerate}
    \item непрерывность должна поддерживаться глобально;
    \item размерностный рост требует несовместимых вариаций, которые нельзя
          представить на одной оси;
    \item на $K_1$ не возникают временные явления;
    \item коллапс при чрезмерных градиентах универсален;
    \item все континуумы $K_1$ должны быть вложены в своё мета-пространство $M_1$.
\end{enumerate}

% ================================================================
\subsubsection{Эксперименты для \texorpdfstring{$K_1$}{K_1}}

Косвенные тесты:
\begin{itemize}
    \item универсальность поведения градиентного коллапса в 1D в физике,
          химии и биологических волокнах,
    \item наблюдение порогового поведения при резких градиентах,
    \item обнаружение вынужденного размерностного роста в системах,
          где одна ось становится недостаточной.
\end{itemize}

Примеры:
\begin{itemize}
    \item схлопывание полимерных цепочек под натяжением,
    \item разрушение мемристорного канала,
    \item нестабильность 1D-каналов реакция-диффузии.
\end{itemize}

% ================================================================
\subsubsection{Смерть \texorpdfstring{$K_1$}{K_1}}

$K_1$ коллапсирует, когда:
\[
\Omega(K_1) = \varnothing.
\]

Это происходит, если:
\begin{itemize}
    \item непрерывность нарушена,
    \item градиенты превышают пороги,
    \item ограничения в $\mathcal{L}$ (переносятся из $K_0$) становятся
          несовместимыми.
\end{itemize}

После коллапса переход к $K_2$ невозможен.

% ================================================================
\subsubsection{Фальсифицируемость \texorpdfstring{$K_1$}{K_1}}

Предсказания, которые могут быть эмпирически опровергнуты:
\begin{enumerate}
    \item универсальное существование порогов непрерывности,
    \item отсутствие нетривиальных циклов,
    \item отсутствие возникновенного времени,
    \item требование монотонной связной оси,
    \item размерностный переход только при несовместимых различиях.
\end{enumerate}

Опровержение любого из этих предсказаний ставит под вопрос описание $K_1$.

% ================================================================
\subsubsection{Ветвление / Онтологическая позиция \texorpdfstring{$K_1$}{K_1}}

$K_1$ располагается непосредственно над $K_0$ в онтологическом дереве:
\[
K_0 \to K_1 \to K_2.
\]

Ветвление в $K_1$ соответствует несовместимым допустимым многообразиям,
имеющим одни и те же мета-аксиомы, но отличающимся структурными
ограничениями.

Такие ветви остаются допустимыми, пока они находятся в $\Omega(K_1)$.

% ================================================================
\subsubsection{Связь с M-пространствами}

$K_1$ вложен в своё мета-пространство:
\[
K_1 \subseteq M_1 \subseteq \Omega(K_0).
\]

$M_1$ определяет:
\begin{itemize}
    \item топологию $\tau$ области $X$,
    \item ограничения регулярности,
    \item допустимые диапазоны потенциалов,
    \item пороги допустимости.
\end{itemize}

Переход к $K_2$ происходит, когда:
\[
A_1 \text{ не может представить все возникающие различия.}
\]
